% =============================================================================
% APPENDIX E: FORMAL PROOFS
%
% Full proofs for all lemmas, propositions, and corollaries.
% The main text retains proof sketches for readability;
% this appendix provides the complete arguments.
% =============================================================================

\section{Formal Proofs}
\label{app:proofs}

This appendix provides full proofs for the results stated in
Sections~\ref{sec:self_balancing}--\ref{sec:multi_agent}. Each proof proceeds
from the measure-space axioms restated below and the definitions of
Section~\ref{sec:definitions}.

% ---------------------------------------------------------------------------
% E.1: Measure Space Axioms
% ---------------------------------------------------------------------------

\subsection{Measure Space Axioms}
\label{app:axioms}

The following properties of $\vol(\cdot)$ are assumed throughout. They
restate Definition~\ref{def:social_objective} and the standing assumptions of
Section~\ref{sec:definitions} in a form convenient for the proofs.

\begin{enumerate}
\item[\textbf{(M1)}] \emph{Non-negativity.} $\vol(A) \geq 0$ for every measurable $A \subseteq \G$.
\item[\textbf{(M2)}] \emph{Null empty set.} $\vol(\emptyset) = 0$.
\item[\textbf{(M3)}] \emph{Monotonicity.} $A \subseteq B \implies \vol(A) \leq \vol(B)$.
\item[\textbf{(M4)}] \emph{Finite additivity.} If $A \cap B = \emptyset$, then $\vol(A \cup B) = \vol(A) + \vol(B)$.
\item[\textbf{(M5)}] \emph{Non-triviality.} For every capability $g \in \G$, the singleton $\{g\}$ has $\vol(\{g\}) > 0$. (Standing assumption from Section~\ref{sec:definitions}.)
\item[\textbf{(M6)}] \emph{Superadditivity under independence.} When agents' individual capability sets are non-overlapping and non-interacting, $V(G) \geq \sum_k \vol(\Gind{k})$, with strict inequality whenever $\Gcoop \neq \emptyset$. (Definition~\ref{def:social_objective}.)
\end{enumerate}

Axioms (M1)--(M4) are standard properties of a finite measure on a
$\sigma$-algebra. (M5) is non-standard: it requires that the measure resolves
individual capabilities, which constrains the granularity of $\G$. (M6) is the
superadditivity assumption from Definition~\ref{def:social_objective}.

% ---------------------------------------------------------------------------
% E.2: Proof of Lemma 2 (Coercion)
% ---------------------------------------------------------------------------

\subsection{Proof of Lemma~\ref{lem:coercion} (Coercive Actions)}
\label{app:proof_coercion}

\begin{proof}
Let $\pi$ restrict agent $a_j$'s capabilities from $\Gind{j}$ to
${\Gind{j}}' \subset \Gind{j}$. Define the restricted joint space as
$G' = \bigl(\bigcup_{i \neq j} \Gind{i} \cup {\Gind{j}}'\bigr) \cup
{\Gcoop}'$, where ${\Gcoop}' = \{g \in \Gcoop : g \text{ remains achievable
under } {\Gind{j}}'\}$. Since ${\Gind{j}}' \subset \Gind{j}$, any
cooperative capability that required a capability in $\Gind{j} \setminus
{\Gind{j}}'$ is no longer achievable, so ${\Gcoop}' \subseteq \Gcoop$.
Combined with ${\Gind{j}}' \subset \Gind{j}$, we have $G' \subseteq G$.
By (M3), $\vol(G') \leq \vol(G)$, giving $\Delta \vol(G) \leq 0$.

For strictness: suppose there exists $g \in \Gind{j} \setminus {\Gind{j}}'$
with $g \notin \bigcup_{i \neq j} \Gind{i}$. Then $g \in G \setminus G'$,
so $G' \subsetneq G$. Write $G = G' \cup (G \setminus G')$ where the union
is disjoint. By (M4), $\vol(G) = \vol(G') + \vol(G \setminus G')$. Since
$\{g\} \subseteq G \setminus G'$, by (M3) and (M5),
$\vol(G \setminus G') \geq \vol(\{g\}) > 0$, giving
$\vol(G) > \vol(G')$.
\end{proof}

% ---------------------------------------------------------------------------
% E.3: Proof of Lemma 3 (Rigidity)
% ---------------------------------------------------------------------------

\subsection{Proof of Lemma~\ref{lem:rigidity} (Self-Imposed Rigidity)}
\label{app:proof_rigidity}

\begin{proof}
Since $A \setminus R \subseteq A$, the feasible set of the restricted problem
is a subset of the feasible set of the unrestricted problem. Therefore
\[
\max_{\pi \in A \setminus R} \mathbb{E}[\Delta \vol(G) \mid \pi]
\;\leq\; \max_{\pi \in A} \mathbb{E}[\Delta \vol(G) \mid \pi].
\]
For strictness: let $\pi^* = \arg\max_{\pi \in A} \mathbb{E}[\Delta \vol(G)
\mid \pi]$. If $\pi^* \in R$ (the unrestricted optimum is forbidden), then
$\pi^*$ is not feasible in $A \setminus R$, so
$\max_{\pi \in A \setminus R} \mathbb{E}[\Delta \vol(G) \mid \pi]$ is attained
at some $\pi' \neq \pi^*$ with
$\mathbb{E}[\Delta \vol(G) \mid \pi'] \leq \mathbb{E}[\Delta \vol(G) \mid
\pi^*]$. If $\pi^*$ is the unique maximizer, the inequality is strict. If
$\pi^*$ is not unique, strictness holds whenever every co-maximizer also
lies in $R$.
\end{proof}

% ---------------------------------------------------------------------------
% E.4: Proof of Corollary 1.2 (Rigid Rules)
% ---------------------------------------------------------------------------

\subsection{Proof of Corollary~\ref{cor:rigid_rules} (Rigid Rules are Anti-Maximizing)}
\label{app:proof_rigid_rules}

\begin{proof}
By Lemma~\ref{lem:rigidity}, strict suboptimality holds whenever
$\pi^* \in R$. It remains to show that such states exist.

Call $R$ \emph{non-trivial} if there exists a reachable state $s$ and an
action $\pi_0 \in R$ with $\mathbb{E}[\Delta \vol(G) \mid \pi_0, s] > 0$:
at least one forbidden action is $\vol(G)$-expanding in some state. For any
non-trivial $R$, let $s_0$ be the state where $\pi_0$ is expanding. If
$\pi_0$ also maximizes $\mathbb{E}[\Delta \vol(G) \mid \cdot\,, s_0]$ over
$A$, we are done. If not, we appeal to an adversarial construction: an adversary who knows the
structure of $R$ can construct a state (via deadline pressure) in which every
action in $A \setminus R$ is $\vol(G)$-contracting while a forbidden action
(e.g., deception, calculated risk) is $\vol(G)$-expanding. In that constructed
state, $\pi^* \in R$ and Lemma~\ref{lem:rigidity} gives strict inequality.
The concrete instance following Corollary~\ref{cor:rigid_rules} provides one
such construction; the argument generalizes to any $R$ whose constraints are
known to the adversary, since knowledge of the permitted region $A \setminus R$
suffices to design scenarios that trap the actor within it.
\end{proof}

% ---------------------------------------------------------------------------
% E.5: Proof of Lemma 1 (Elimination)
% ---------------------------------------------------------------------------

\subsection{Proof of Lemma~\ref{lem:destruction} (Net Effect of Elimination)}
\label{app:proof_elimination}

\begin{proof}
\emph{Setup.} Partition the joint goal space $G$ into four disjoint regions
relative to agent $a_k$:
\begin{align*}
U_k &= \Gind{k} \setminus \bigcup_{j \neq k} \Gind{j}
  && \text{($a_k$'s unique individual capabilities)} \\
S_k &= \Gind{k} \cap \bigcup_{j \neq k} \Gind{j}
  && \text{(shared individual capabilities)} \\
K_k &= G^{\mathrm{coop}}_{k}
  && \text{(cooperative capabilities requiring $a_k$)} \\
R   &= G \setminus (U_k \cup S_k \cup K_k)
  && \text{(everything else: other agents' unique caps + coop not involving $a_k$)}
\end{align*}
By construction these four sets are disjoint and
$G = U_k \cup S_k \cup K_k \cup R$.
By (M4), $\vol(G) = \vol(U_k) + \vol(S_k) + \vol(K_k) + \vol(R)$.

\emph{After elimination.} Removing $a_k$ produces $G'$. The shared
capabilities $S_k$ survive (other agents possess them). The remainder $R$
is unaffected. The unique capabilities $U_k$ are lost. The cooperative
capabilities $K_k$ are lost. Additionally, removing $a_k$ may unlock
\emph{suppression-released} capabilities
$G^{\mathrm{coop}+}_{k} \subseteq \G \setminus G$: cooperative capabilities
among the remaining agents that were not achievable while $a_k$ was present
(e.g., because $a_k$'s interference prevented the required coordination).
By definition, $G^{\mathrm{coop}+}_{k} \cap G = \emptyset$, since these
capabilities were not in $G$ before removal.

The post-elimination space is therefore
$G' = S_k \cup R \cup G^{\mathrm{coop}+}_{k}$, and since
$G^{\mathrm{coop}+}_{k}$ is disjoint from $S_k \cup R$ (it was not in $G$),
\[
\vol(G') = \vol(S_k) + \vol(R) + \vol(G^{\mathrm{coop}+}_{k}).
\]

\emph{Volume change.}
\begin{align*}
\Delta \vol(G) &= \vol(G') - \vol(G) \\
  &= \bigl[\vol(S_k) + \vol(R) + \vol(G^{\mathrm{coop}+}_{k})\bigr]
   - \bigl[\vol(U_k) + \vol(S_k) + \vol(K_k) + \vol(R)\bigr] \\
  &= \vol(G^{\mathrm{coop}+}_{k}) - \vol(U_k) - \vol(K_k)
\end{align*}
which is the claimed Equation~\eqref{eq:destruction_bound}.

\emph{Sign.} Elimination is net-contracting ($\Delta \vol(G) < 0$) when
$\vol(U_k) + \vol(K_k) > \vol(G^{\mathrm{coop}+}_{k})$. Under (M5), if
$a_k$ has at least one unique capability, $\vol(U_k) > 0$; if $a_k$
participates in at least one coalition, $\vol(K_k) > 0$. Since
$G^{\mathrm{coop}+}_{k}$ is typically empty or small (most agents do not
suppress cooperation by their mere presence), the condition holds for generic
agents. For scorpions that actively suppress cooperation,
$\vol(G^{\mathrm{coop}+}_{k})$ can dominate, making elimination
net-expanding.
\end{proof}

% ---------------------------------------------------------------------------
% E.6: Proof of Corollary 1.1 (Elimination Cost)
% ---------------------------------------------------------------------------

\subsection{Proof of Corollary~\ref{cor:elimination} (Elimination is Almost Always Anti-Maximizing)}
\label{app:proof_elimination_cost}

\begin{proof}
\emph{Part 1: Positive unique contribution.} Suppose $a_k$ has at least one
capability $g \in \Gind{k}$ with $g \notin \Gind{j}$ for all $j \neq k$.
Then $\{g\} \subseteq U_k = \Gind{k} \setminus \bigcup_{j \neq k} \Gind{j}$.
By (M3) and (M5), $\vol(U_k) \geq \vol(\{g\}) > 0$.

\emph{Part 2: Cooperative loss grows with interaction partners.} For each
agent $a_j \neq a_k$ that interacts with $a_k$, define the pairwise
cooperative contribution $G^{\mathrm{coop}}_{\{k,j\}}$ as the set of
capabilities achievable by the pair $\{a_k, a_j\}$ but not by either alone
and not by any other pair not involving $a_k$. The sets
$\{G^{\mathrm{coop}}_{\{k,j\}}\}_{j \neq k}$ are not necessarily disjoint
(a capability may require $a_k$ and either of two partners), so we cannot
directly sum their volumes.

However, for each interacting partner $a_j$, assume there exists at least one
capability $g_j \in G^{\mathrm{coop}}_{\{k,j\}}$ that is unique to that pair
(not achievable by $a_k$ with any other single partner). Under (M5),
$\vol(\{g_j\}) > 0$. The set $\{g_j : j \text{ interacts with } a_k\}$
has cardinality equal to the number of interaction partners $m_k$. If the
$\{g_j\}$ are distinct, then by (M4) applied to the disjoint singletons,
$\vol(G^{\mathrm{coop}}_{k}) \geq \sum_{j} \vol(\{g_j\}) > 0$, and this
bound grows linearly in $m_k$.

The combinatorial observation in the main text (that $2^{n-1} - 1$ coalitions
involving $a_k$ exist) motivates the expectation of super-linear growth.
Formalizing this requires a \emph{diversity condition}: each coalition
$S \ni a_k$ with $|S| \geq 2$ produces at least one capability not achievable
by any proper sub-coalition of $S$ containing $a_k$. Under this condition,
the number of unique coalition-level capabilities grows with the number of
coalitions, but the precise scaling depends on the overlap structure among
coalitions and is not characterized here. The linear lower bound in $m_k$
holds without the diversity condition.
\end{proof}

% ---------------------------------------------------------------------------
% E.7: Proof of Lemma 4 (Cooperative Expansion)
% ---------------------------------------------------------------------------

\subsection{Proof of Lemma~\ref{lem:cooperation} (Cooperative Expansion)}
\label{app:proof_cooperation}

\begin{proof}
\emph{Part (i).} By definition, $G^{\mathrm{coop}}_{ij}(t)$ contains
capabilities requiring the simultaneous participation of both $a_i$ and
$a_j$. A unilateral action $\pi$ by $a_i$ alone does not involve $a_j$'s
participation. No capability in $G^{\mathrm{coop}}_{ij}$ can be created or
destroyed by $\pi$, since its defining condition (joint participation) is
not met. Therefore $G^{\mathrm{coop}}_{ij}(t+1) = G^{\mathrm{coop}}_{ij}(t)$
under $\pi$, giving $\Delta \vol(G^{\mathrm{coop}}_{ij} \mid \pi) = 0$.
By symmetry the same holds for unilateral actions by $a_j$.

\emph{Part (ii).} Call $a_i$ and $a_j$ \emph{complementary at time $t$} if
there exists a joint action profile $(\pi_i, \pi_j) \in A_i \times A_j$ that
produces a capability $g \notin G(t)$: a capability not achievable under the
current coordination infrastructure. Suppose $a_i$ and $a_j$ are
complementary and have not yet established the coordination mechanism that
enables $g$ (e.g., a communication channel, shared protocol, or mutual
commitment). Then $g \notin G(t)$. The joint action $\pi_{ij}$ that
establishes the mechanism makes $g$ achievable, so
$g \in G^{\mathrm{coop}}_{ij}(t+1)$ and
$G^{\mathrm{coop}}_{ij}(t+1) \supsetneq G^{\mathrm{coop}}_{ij}(t)$. By (M5),
$\vol(\{g\}) > 0$, so
$\vol(G^{\mathrm{coop}}_{ij}(t+1)) \geq \vol(G^{\mathrm{coop}}_{ij}(t)) +
\vol(\{g\}) > \vol(G^{\mathrm{coop}}_{ij}(t))$,
giving $\Delta \vol(G^{\mathrm{coop}}_{ij} \mid \pi_{ij}) > 0$.
Whether this cooperative gain produces a net expansion in $\vol(G)$ depends on
whether the opportunity cost $C_{ij}$ of coordination (individual capabilities
foregone) exceeds the cooperative gain; part~(iii) addresses this.

\emph{Part (iii).} Let $\Delta^{\mathrm{coop}} =
\mathbb{E}[\Delta \vol(G^{\mathrm{coop}}_{ij})]$ be the expected cooperative
expansion from joint action, and let $C_{ij} \geq 0$ be the opportunity cost:
the $\vol(G)$ value of the best individual actions $a_i$ and $a_j$ forgo in
order to coordinate. The net expected change in $\vol(G)$ from cooperation is
$\Delta^{\mathrm{coop}} - C_{ij}$, which is positive whenever
$\Delta^{\mathrm{coop}} > C_{ij}$. A GFM actor maximizing
$\mathbb{E}[\Delta \vol(G)]$ therefore prefers cooperation to unilateral
action when this condition holds.
\end{proof}

% ---------------------------------------------------------------------------
% E.8: Proof of Proposition 3 (Sign-Correctness)
% ---------------------------------------------------------------------------

\subsection{Proof of Proposition~\ref{prop:sign_correctness} (Sign-Correctness Decomposition)}
\label{app:proof_sign_correctness}

\begin{proof}
Write the true volume change as
$\Delta \vol(G) = \Delta_I + \Delta_C$, where
$\Delta_I = \Delta \vol\bigl(\bigcup_k \Gind{k}\bigr)$ is the
individual-capability component and $\Delta_C = \Delta \vol(\Gcoop)$ is the
cooperative component. The estimator observes
$\hat{\Delta}_I \approx \Delta_I$ (conditions 1--2 control the approximation
error) and $\hat{\Delta}_C$ with error
$e_C = \Delta_C - \hat{\Delta}_C$.

The estimator's output is
$\Deltahat(\pi) = \hat{\Delta}_I + \hat{\Delta}_C
= (\Delta_I - e_I) + (\Delta_C - e_C)
= \Delta \vol(G) - (e_I + e_C)$,
where $e_I = \Delta_I - \hat{\Delta}_I$ is the individual-level error.

Under conditions 1--2, $|e_I|$ is small relative to $|\Delta_I|$ (the actor
observes a representative sample with an accurate trust model). Under
condition~3, $|e_C| < (1 - \epsilon) \cdot |\Delta_I|$.

For sign-correctness we need $\sign(\Deltahat(\pi)) = \sign(\Delta \vol(G))$.
Since $\Deltahat(\pi) = \Delta \vol(G) - (e_I + e_C)$, the sign is preserved
when $|e_I + e_C| < |\Delta \vol(G)|$. By the triangle inequality,
$|e_I + e_C| \leq |e_I| + |e_C|$. Write $|e_I| \leq \delta |\Delta_I|$ for
a sampling error ratio $\delta$ controlled by conditions 1--2, and
$|e_C| < (1-\epsilon) |\Delta_I|$ by condition~3. Two cases:

\emph{Case 1: $\Delta_I$ and $\Delta_C$ have the same sign.} Then
$|\Delta \vol(G)| = |\Delta_I| + |\Delta_C| \geq |\Delta_I|$, and
$|e_I| + |e_C| < (\delta + 1 - \epsilon)|\Delta_I|$, which is less than
$|\Delta_I| \leq |\Delta \vol(G)|$ whenever $\delta < \epsilon$.

\emph{Case 2: $\Delta_I$ and $\Delta_C$ have opposite signs.} Then
$|\Delta \vol(G)| = \bigl||\Delta_I| - |\Delta_C|\bigr|$, and sign
preservation requires $(\delta + 1 - \epsilon)|\Delta_I| <
\bigl||\Delta_I| - |\Delta_C|\bigr|$. This holds when
$|\Delta_C| < (\epsilon - \delta)|\Delta_I|$, a condition strictly stronger
than condition~3 alone. Condition~3 guarantees sign-correctness in this case
only when the cooperative-level \emph{change} (not just the estimation error)
is small relative to the individual-level change.
\end{proof}

% ---------------------------------------------------------------------------
% E.9: Proof of Proposition 1 (Self-Balancing)
% ---------------------------------------------------------------------------

\subsection{Proof of Proposition~\ref{prop:self_balancing} (Self-Balancing Property)}
\label{app:proof_self_balancing}

\begin{proof}
\emph{Parts (a)--(c)} follow directly from the lemmas.

(a) By Lemma~\ref{lem:destruction}, eliminating agent $a_k$ produces
$\Delta \vol(G) = \vol(G^{\mathrm{coop}+}_{k}) - \vol(U_k) - \vol(K_k)$.
When $\vol(U_k) + \vol(K_k) > \vol(G^{\mathrm{coop}+}_{k})$ (the generic
case for non-scorpion agents), destruction is $\vol(G)$-contracting.

(b) By Lemma~\ref{lem:coercion}, restricting $a_j$'s capabilities produces
$G' \subseteq G$ and $\vol(G') \leq \vol(G)$, with strict inequality when the
restriction removes capabilities not covered by other agents.

(c) By Lemma~\ref{lem:rigidity}, constraining the actor to $A \setminus R$
gives $\max_{A \setminus R} \mathbb{E}[\Delta \vol(G)] \leq \max_A
\mathbb{E}[\Delta \vol(G)]$, with strict inequality when $R$ excludes the
unrestricted optimum.

\emph{Part (d): Partial result under additional assumptions.} Parts (a)--(c)
establish that the $\vol(G)$ objective penalizes movement toward both extremes:
excessive destruction contracts $G$ by (a), excessive restriction contracts $G$
by (b), and rigid self-constraint reduces the actor's ability to expand $G$
by (c). To show that these opposing pressures produce an optimum distinct from
both extremes, assume:
\begin{itemize}
\item The action space $A$ is compact.
\item The map $\pi \mapsto \mathbb{E}[\Delta \vol(G) \mid \pi]$ is continuous.
\item There exist actions $\pi_d, \pi_r \in A$ such that $\pi_d$
  is destructive ($\mathbb{E}[\Delta \vol(G) \mid \pi_d] < 0$ by (a)) and
  $\pi_r$ is overrestrictive ($\mathbb{E}[\Delta \vol(G) \mid \pi_r] < 0$
  by (b)).
\item There exists $\pi_0 \in A$ with
  $\mathbb{E}[\Delta \vol(G) \mid \pi_0] \geq 0$.
\end{itemize}
By the extreme value theorem (continuity on a compact set), the maximum of
$\mathbb{E}[\Delta \vol(G) \mid \pi]$ over $A$ is attained at some
$\pi^* \in A$. Since $\pi_0$ achieves a non-negative value,
$\mathbb{E}[\Delta \vol(G) \mid \pi^*] \geq 0$. Since both $\pi_d$ and
$\pi_r$ yield strictly negative values, $\pi^* \neq \pi_d$ and
$\pi^* \neq \pi_r$. The maximizer is therefore distinct from the destructive
and overrestrictive extremes.

This argument establishes \emph{existence} of an optimum that is neither
maximally destructive nor maximally restrictive, under continuity and
compactness. It does not establish \emph{stability} (whether small
perturbations return to the optimum), \emph{uniqueness} (whether there is one
such point or many), or \emph{convergence} (whether the optimization loop of
Section~\ref{subsec:optimization_loop} reaches it). These remain open
problems.
\end{proof}

% ---------------------------------------------------------------------------
% E.10: Proof of Proposition 2 (Scorpion Detection)
% ---------------------------------------------------------------------------

\subsection{Proof of Proposition~\ref{prop:scorpion_detection} (Scorpion Detection)}
\label{app:proof_scorpion}

\begin{proof}
\emph{Channel (a): Contraction attribution.}
Let $a_j$ be a scorpion whose actions persistently contract $\vol(G)$. At
each time step $t$, the GFM actor observes the capability changes
$\Delta \vol(\Gind{k} \mid t)$ for each $a_k \in \mathcal{O}$ (the
observation set). When $a_j$ acts, the affected agents report
$R_k(t) \in \{-1, 0, +1\}$. If $a_j$'s actions persistently reduce
capabilities, the trust-weighted aggregate
$\Deltahat_j(t) = \sum_{k \in \mathcal{O}} T_k \cdot R_k(t)$ is
persistently negative in time steps where $a_j$ acts.

Formally, define the running average contraction signal attributed to $a_j$:
\[
\bar{\Delta}_j(\tau) = \frac{1}{\tau} \sum_{t=1}^{\tau}
  \Deltahat_j(t) \cdot \mathbf{1}[a_j \text{ acted at } t].
\]
If $a_j$'s actions produce expected contraction $\mu_j < 0$, the observation
noise has finite variance $\sigma^2$, and the observation windows
(time steps where $a_j$ acts) yield conditionally independent signals given
$a_j$'s strategy, then by the strong law of large numbers,
$\bar{\Delta}_j(\tau) \to \mu_j < 0$ almost surely as $\tau \to \infty$.
After $\tau^* = O(\sigma^2 / \mu_j^2)$ observations, the signal exceeds any
fixed detection threshold with high probability. The independence condition
is non-trivial: successive observations of a scorpion's effects may be
correlated when the same agents are repeatedly affected. Under weaker mixing
conditions the convergence still holds but the rate degrades.

This is a \emph{correlational} detection: the signal identifies that negative
capability changes co-occur with $a_j$'s actions. Causal attribution
(separating $a_j$'s contribution from simultaneous causes) requires
additional structure not provided by the framework.

\emph{Channel (b): Deception detection.}
If $a_j$ reports capability changes that diverge from independently observed
outcomes, the prediction residual $r_j(t) = R_j(t) - \hat{R}_j(t \mid
M_j(G))$ is nonzero. The cumulative inconsistency
$\sigma_j^2(t) = \alpha \cdot \sigma_j^2(t-1) + (1-\alpha) \cdot r_j(t)^2$
(Equation~\eqref{eq:cumulative_consistency}) is an exponentially weighted
moving average (EWMA) of squared residuals.

Under stationarity of $a_j$'s deceptive behavior, $\sigma_j^2(t)$ converges
to $\mathbb{E}[r_j^2]$ (Appendix~\ref{app:trust_update},
Equation~\eqref{eq:convergence}). For a deceptive scorpion,
$\mathbb{E}[r_j^2] > 0$, so
$T_j \to T_j^* = 1/(1 + \beta \cdot \mathbb{E}[r_j^2]) < 1$
(Equation~\eqref{eq:trust_factor}). The convergence rate is controlled by
the decay parameter $\alpha$: after approximately $1/(1-\alpha)$ observations,
the prior's influence is negligible and $T_j$ reflects the agent's actual
prediction consistency.

For an honest (non-deceptive) cooperative agent, $\mathbb{E}[r_j^2]$ is
small (driven by observation noise), so $T_j^* \approx 1$. The gap
$T_j^*(\text{deceptive}) < T_j^*(\text{honest})$ provides a detection
threshold that improves with observation count.

\emph{Convergence caveat.} Both channels require stationarity of the
scorpion's behavior. A scorpion that alternates between cooperative and
contracting phases, or that changes its strategy in response to detection
pressure, can maintain $\bar{\Delta}_j \approx 0$ and $T_j \approx 1$
indefinitely. Detection of non-stationary adversaries requires adaptive
methods beyond the EWMA framework and remains an open problem.
\end{proof}
